\documentclass[11pt]{article}

% ---------- Packages ----------
\usepackage[margin=1in]{geometry}
\usepackage{amsmath, amssymb, amsthm}
\usepackage{mathtools}
\usepackage{enumitem}
\usepackage{graphicx}
\usepackage{booktabs}
\usepackage{array}
\usepackage{hyperref}
\usepackage{xcolor}
\usepackage{fancyhdr}
\usepackage{lastpage}
\usepackage{listings}
\usepackage{courier}

% ---------- Hyperref ----------
\hypersetup{
    colorlinks=true,
    linkcolor=blue,
    urlcolor=blue
}

% ---------- Header/Footer ----------
\pagestyle{fancy}
\fancyhf{}

\lhead{MIT 6.1200J}
\chead{Problem Set 10}
\rhead{Nicolas Iskandar}

\cfoot{\thepage\ /\ \pageref{LastPage}}

% ---------- Custom Commands ----------
\newcommand{\R}{\mathbb{R}}
\newcommand{\N}{\mathbb{N}}
\newcommand{\Z}{\mathbb{Z}}

\newcommand{\answer}{
    \vspace{0.5em}
    \noindent\textbf{Answer.}
}

% ---------- Code Style ----------
\lstset{
    basicstyle=\ttfamily\small,
    breaklines=true,
    frame=single,
    columns=fullflexible,
    keepspaces=true
}

% ---------- Title ----------
\title{
    \vspace{-2em}
    MIT 6.1200J \\
    Problem Set 10
}
\author{Nicolas Iskandar}
\date{}

% ==================================================
\begin{document}

\maketitle

\section*{Problem 1: Random Pixels}
\begin{enumerate}[label=\textbf{(\alph*)}]
    \item We want to find $\Pr[A]$, the probability that the 4 cells around vertex $a$ (cells 2, 3, 6, 7) all share the same color. 
    There are 4 possible colors they could all share. For any specific color, the probability that all 4 independent cells are that color is $(1/4)^4$. 
    Because these are 4 mutually exclusive cases, we add their probabilities:
    $$\Pr[A] = 4 \times \left(\frac{1}{4}\right)^4 = \frac{1}{4^3} = \frac{1}{64}$$

    \item To prove $A, B,$ and $C$ are mutually independent, we must show that the probability of their intersection equals the product of their individual probabilities for all possible combinations. We already established $\Pr[A] = \Pr[B] = \Pr[C] = 1/64$.
    
    \begin{itemize}
        \item \textbf{Pairwise $A \cap B$:} The set of cells for $A$ is $S_A = \{2, 3, 6, 7\}$ and for $B$ is $S_B = \{7, 8, 11, 12\}$. They share exactly one cell (cell 7). For $A$ and $B$ to both occur, all cells in $S_A$ must be color $c_1$ and all cells in $S_B$ must be color $c_2$. Because they overlap at cell 7, $c_1$ must equal $c_2$. Thus, all $4+4-1=7$ unique cells in $S_A \cup S_B$ must be the same color. 
        $$\Pr[A \cap B] = 4 \times \left(\frac{1}{4}\right)^7 = \frac{1}{4^6} = \left(\frac{1}{64}\right)^2 = \Pr[A]\Pr[B]$$
        
        \item \textbf{Pairwise $B \cap C$:} $S_B = \{7, 8, 11, 12\}$ and $S_C = \{10, 11, 14, 15\}$. They share exactly cell 11. By the exact same logic as above, all 7 cells in their union must be the same color.
        $$\Pr[B \cap C] = 4 \times \left(\frac{1}{4}\right)^7 = \frac{1}{4^6} = \Pr[B]\Pr[C]$$

        \item \textbf{Pairwise $A \cap C$:} $S_A$ and $S_C$ share no cells. Thus, the 4 cells in $S_A$ can be a color $c_1$ and the 4 cells in $S_C$ can independently be a color $c_2$. 
        $$\Pr[A \cap C] = \left( 4 \times \left(\frac{1}{4}\right)^4 \right) \times \left( 4 \times \left(\frac{1}{4}\right)^4 \right) = \frac{1}{64} \times \frac{1}{64} = \Pr[A]\Pr[C]$$

        \item \textbf{Three-way $A \cap B \cap C$:} We examine the union $S_A \cup S_B \cup S_C$. As established, $A \cap B$ forces $c_A = c_B$, and $B \cap C$ forces $c_B = c_C$. Therefore, all cells across the three sets must share a single color. There are $4 + 3 + 3 = 10$ unique cells.
        $$\Pr[A \cap B \cap C] = 4 \times \left(\frac{1}{4}\right)^{10} = \frac{1}{4^9} = \left(\frac{1}{4^3}\right)^3 = \Pr[A]\Pr[B]\Pr[C]$$
    \end{itemize}
    Since all 4 conditions hold, $A, B,$ and $C$ are mutually independent.

    \item To see if $A, B, C,$ and $D$ are mutually independent, we check $\Pr[A \cap B \cap C \cap D]$.
    The set of cells for $D$ is $S_D = \{5, 6, 9, 10\}$. 
    If all four events occur:
    \begin{itemize}
        \item $A \cap B$ forces $c_A = c_B$ (overlap at 7)
        \item $B \cap C$ forces $c_B = c_C$ (overlap at 11)
        \item $C \cap D$ forces $c_C = c_D$ (overlap at 10)
    \end{itemize}
    This chain forces all 4 events to share the same color. The union of all four sets $S_A \cup S_B \cup S_C \cup S_D$ contains 12 unique cells (everything except the 4 corners: 1, 4, 13, 16).
    The probability that all 12 cells share the same color is:
    $$\Pr[A \cap B \cap C \cap D] = 4 \times \left(\frac{1}{4}\right)^{12} = \frac{1}{4^{11}}$$
    If they were mutually independent, the product of their probabilities would be:
    $$\Pr[A]\Pr[B]\Pr[C]\Pr[D] = \left(\frac{1}{4^3}\right)^4 = \frac{1}{4^{12}}$$
    Because $1/4^{11} \neq 1/4^{12}$, the four events are \textbf{not} mutually independent.
\end{enumerate}

\vspace{0.5cm}

\section*{Problem 2: Median vs Maximum}
The total number of ways to draw 3 Raboots from 100 without replacement is $\binom{100}{3}$.

\begin{enumerate}[label=\textbf{(\alph*)}]
    \item \textbf{Values of $PMF_D(k)$ and $PMF_X(k)$:}
    
    For the median $D = k$, we need exactly one draw to be less than $k$, one draw to be exactly $k$, and one draw to be greater than $k$. 
    There are $k-1$ choices for the smaller number and $100-k$ choices for the larger number.
    $$ PMF_D(k) = \frac{\binom{k-1}{1} \binom{100-k}{1}}{\binom{100}{3}} = \frac{(k-1)(100-k)}{\binom{100}{3}} $$
    
    For the maximum $X = k$, we need exactly two draws to be less than $k$, and one draw to be exactly $k$.
    There are $\binom{k-1}{2}$ ways to pick the two smaller numbers.
    $$ PMF_X(k) = \frac{\binom{k-1}{2}}{\binom{100}{3}} $$

    \item \textbf{Value of $CDF_D(k)$:}
    The Cumulative Distribution Function $CDF_D(k)$ is the probability that the median is $\le k$. This occurs if and only if \textbf{at least two} of the three drawn numbers are $\le k$. We split this into two disjoint cases:
    \begin{itemize}
        \item \textbf{Exactly two numbers $\le k$:} We choose 2 numbers from the $k$ available, and 1 number from the $100-k$ numbers strictly greater than $k$. Ways = $\binom{k}{2}\binom{100-k}{1}$.
        \item \textbf{Exactly three numbers $\le k$:} We choose all 3 numbers from the $k$ available. Ways = $\binom{k}{3}$.
    \end{itemize}
    $$ CDF_D(k) = \frac{\binom{k}{2}(100-k) + \binom{k}{3}}{\binom{100}{3}} $$

    \item \textbf{$\Pr[D = 10 \text{ AND } X = 5]$:}
    The median of a set of numbers can never be strictly greater than the maximum of that same set. Since $10 > 5$, this event is impossible.
    $$ \Pr[D = 10 \cap X = 5] = 0 $$

    \item \textbf{Are $D$ and $X$ independent?}
    No. Two random variables are independent if $\Pr[D=a \cap X=b] = \Pr[D=a]\Pr[X=b]$ for all possible values. 
    From part (c), $\Pr[D=10 \cap X=5] = 0$. However, $\Pr[D=10] > 0$ and $\Pr[X=5] > 0$, making their product strictly positive. Since $0 \neq \text{positive}$, they are dependent.

    \item \textbf{Support of $PMF_{X-D}(k)$:}
    $X - D$ represents the difference between the maximum and the median. Since all three drawn numbers are distinct integers, the maximum must be strictly greater than the median. Thus, the minimum difference is $X - D \ge 1$.
    The maximum possible value for $X$ is 100, and the smallest possible value for $D$ when $X=100$ is 2 (e.g., drawing 1, 2, 100). The difference here is $100 - 2 = 98$.
    Since the draws are integers, the difference must be an integer. Thus, $PMF_{X-D}(k) > 0$ for all integers $k$ such that:
    $$ k \in \{1, 2, 3, \dots, 98\} $$
\end{enumerate}

\vspace{0.5cm}

\section*{Problem 3: Probabilistic Method}
\begin{enumerate}[label=\textbf{(\alph*)}]
    \item A single clause $C_j$ has $k$ literals. A literal is False with probability $1/2$. Because the variables are distinct and coin flips are independent, the probability that \textit{all} $k$ literals in the clause evaluate to False is $(1/2)^k$. 
    The probability that the clause is satisfied (at least one True literal) is the complement:
    $$ \Pr[C_j \text{ is satisfied}] = 1 - \frac{1}{2^k} $$

    \item Let $I_j$ be the indicator random variable for the event that clause $C_j$ is \textbf{unsatisfied}. 
    $$ I_j = \begin{cases} 1 & \text{if } C_j \text{ is unsatisfied} \\ 0 & \text{if } C_j \text{ is satisfied} \end{cases} $$
    The total number of unsatisfied clauses $U$ is simply the sum of these indicators:
    $$ U = \sum_{j=1}^m I_j $$

    \item By linearity of expectation, the expected value of a sum is the sum of expected values:
    $$ \text{Ex}[U] = \text{Ex}\left[ \sum_{j=1}^m I_j \right] = \sum_{j=1}^m \text{Ex}[I_j] $$
    The expected value of an indicator variable is simply the probability of its event. From part (a), the probability a clause is unsatisfied is $1/2^k$.
    $$ \text{Ex}[U] = \sum_{j=1}^m \frac{1}{2^k} = \frac{m}{2^k} $$

    \item If $k=5$ and $m=31$, we plug these into our expectation formula:
    $$ \text{Ex}[U] = \frac{31}{2^5} = \frac{31}{32} $$
    Since $U$ is a count of unsatisfied clauses, its possible values are non-negative integers $\{0, 1, 2, \dots, 31\}$. 
    Because the expected (average) value of $U$ is strictly less than 1 ($\frac{31}{32} < 1$), it is mathematically impossible for every single outcome of $U$ to be $\ge 1$. There must exist at least one outcome in the probability space where $U < 1$. Since $U$ is an integer, this outcome must be $U = 0$. 
    Therefore, $\Pr[U=0] > 0$.

    \item From part (d), because $\Pr[U=0] > 0$, the event $U=0$ has a strictly positive chance of occurring. The event $U=0$ exactly describes a scenario where there are 0 unsatisfied clauses, meaning every single clause is satisfied. 
    Because the probability is non-zero, such an assignment must logically exist.
\end{enumerate}

\vspace{0.5cm}

\section*{Problem 4: Teeter Dance}
\begin{enumerate}[label=\textbf{(\alph*)}]
    \item \textbf{PMF Distributions:}
    
    \begin{tabular}{|l|c|c|c|c|c|c|c|c|c|}
    \hline
    \textbf{location} & \textbf{-4} & \textbf{-3} & \textbf{-2} & \textbf{-1} & \textbf{0} & \textbf{1} & \textbf{2} & \textbf{3} & \textbf{4} \\
    \hline
    initially & 0 & 0 & 0 & 0 & 1 & 0 & 0 & 0 & 0 \\
    after 1 step & 0 & 0 & 0 & 1/2 & 0 & 1/2 & 0 & 0 & 0 \\
    after 2 steps & 0 & 0 & 1/4 & 0 & 2/4 & 0 & 1/4 & 0 & 0 \\
    after 3 steps & 0 & 1/8 & 0 & 3/8 & 0 & 3/8 & 0 & 1/8 & 0 \\
    after 4 steps & 1/16 & 0 & 4/16 & 0 & 6/16 & 0 & 4/16 & 0 & 1/16 \\
    \hline
    \end{tabular}

    \item \textbf{Regaining Bearings:}
    \begin{enumerate}[label=(\roman*)]
        \item Let Spinda step East exactly $i$ times. Since the walk has length $t$, Spinda must step West exactly $t-i$ times. 
        Each East step changes location by $+1$, and each West step by $-1$.
        $$ s = (i)(1) + (t-i)(-1) = i - t + i = 2i - t $$
        
        \item A specific walk of length $t$ is fully defined by which $i$ steps out of the $t$ total steps were taken to the East. This is choosing $i$ items from $t$:
        $$ \text{Number of walks} = \binom{t}{i} $$
        
        \item Because each step (East or West) is independent and has probability $1/2$, the probability of any one specific walk sequence of length $t$ is exactly $(1/2)^t$. 
        Since there are $\binom{t}{i}$ distinct walks that end at location $s$, the total probability is:
        $$ \Pr[L_t = s] = \binom{t}{i} \left(\frac{1}{2}\right)^t $$
        
        \item Let $B_t = (L_t + t)/2$. 
        Substitute the expression for location $L_t = 2i - t$ into the equation:
        $$ B_t = \frac{(2i - t) + t}{2} = \frac{2i}{2} = i $$
        We see that the random variable $B_t$ exactly counts $i$, the number of East steps taken in $t$ trials. 
        Since Spinda takes $t$ independent steps (trials) and each has a $1/2$ chance of being East (success), the number of successes $i$ strictly follows a Binomial distribution.
        Thus, $B_t \sim Bin_t(1/2)$.
    \end{enumerate}
\end{enumerate}

\end{document}